I checked whether there was a standard polytope lemma that would make this shorter; for context, Gao notes that in Grünbaum’s sense every plane polygon with more than three edges is homothetically decomposable, but I will not use that here. ([math.clemson.edu](https://www.math.clemson.edu/~sgao/papers/polytope_irr.pdf))

Here is a complete proof.

Let the indices be read modulo \(n\).  
Write
\[
p_i:=s_i+t_i \qquad (1\le i\le n).
\]
These are the \(n\) lonely points coming from Subproblems \(1\)–\(4\). Since
\[
\bigl|\{p:\sigma(p)\neq 0\}\bigr|=n,
\]
they are in fact **all** the points with nonzero \(\sigma\).

Also, as noted in the hint, going once around the boundary gives
\[
\sigma(t_{i+1})=(-1)^{m_i}\sigma(t_i),
\]
hence
\[
(-1)^{m_1+\cdots+m_n}=1,
\]
so \(\sum_i m_i\) is even. I record this, although the contradiction below will come from a more local argument.

## Step \(1\): choose a minimal ear

For each \(i\), define
\[
a_i=s_{i+1}-s_i,
\]
\[
\Delta_i:=\operatorname{conv}(s_{i-1},s_i,s_{i+1}),
\]
and
\[
r_i:=s_i+a_i-a_{i-1}=s_{i-1}+s_{i+1}-s_i.
\]
Let
\[
K_i:=s_i+[0,1]a_i+[0,1](-a_{i-1})
=\operatorname{conv}(s_{i-1},s_i,s_{i+1},r_i).
\]
Thus \(K_i\) is the parallelogram built on the two edges of \(P\) adjacent to \(s_i\).

Choose \(i\) so that \(\operatorname{area}(\Delta_i)\) is minimal.

### Claim
\[
K_i\cap S=\{s_{i-1},s_i,s_{i+1}\}.
\]

### Proof of the claim
Assume not. Then some vertex of \(P\) distinct from \(s_{i-1},s_i,s_{i+1}\) lies in \(K_i\).

Apply an affine map sending
\[
s_i\mapsto (0,0),\qquad s_{i+1}\mapsto (1,0),\qquad s_{i-1}\mapsto (0,1).
\]
In these coordinates,
\[
K_i=[0,1]^2,
\]
\[
\Delta_i=\{(x,y):x\ge 0,\ y\ge 0,\ x+y\le 1\},
\]
and the complementary triangle in \(K_i\) is
\[
T_i:=\operatorname{conv}\bigl((1,0),(1,1),(0,1)\bigr).
\]

Because \(P\) is convex, the boundary arc of \(P\) from \(s_{i+1}\) to \(s_{i-1}\) not containing \(s_i\) is a convex chain; its intersection with the convex triangle \(T_i\) is therefore a connected subchain. Since we assumed there is another vertex of \(P\) in \(K_i\), and no other vertex can lie in \(\Delta_i\), this subchain contains an interior vertex \(v\) of \(P\). Let \(u\) and \(w\) be its two neighbors on \(P\). Then
\[
u,v,w\in T_i.
\]
Hence the ear triangle at \(v\),
\[
\operatorname{conv}(u,v,w),
\]
is contained in \(T_i\). Therefore
\[
\operatorname{area}\bigl(\operatorname{conv}(u,v,w)\bigr)
\le \operatorname{area}(T_i)
= \operatorname{area}(\Delta_i).
\]

If equality held, then \(\operatorname{conv}(u,v,w)=T_i\), so
\[
\{u,v,w\}=\{s_{i+1},r_i,s_{i-1}\}.
\]
That forces the opposite boundary arc from \(s_{i+1}\) to \(s_{i-1}\) to consist only of \(s_{i+1},r_i,s_{i-1}\). Thus
\[
P=\operatorname{conv}(s_{i-1},s_i,s_{i+1},r_i),
\]
which is a parallelogram. But \(P\) has at least three edge-direction classes, so \(P\) is not a parallelogram.

So in fact
\[
\operatorname{area}\bigl(\operatorname{conv}(u,v,w)\bigr)
< \operatorname{area}(\Delta_i),
\]
contradicting the minimal choice of \(i\).

This proves the claim.

---

## Step \(2\): build an extra noncanceling point

By the Subproblem \(4\) structure, the \(t_i\) are consecutive vertices of \(Q:=\operatorname{conv}(T)\), and
\[
t_{i+1}=t_i+m_i a_i
\]
with \(m_i\ge 1\). Hence
\[
t_i-a_{i-1}\in F_{i-1},\qquad t_i+a_i\in F_i.
\]
Set
\[
u:=t_i-a_{i-1},\qquad v:=t_i+a_i.
\]
Since signs alternate along each arithmetic progression,
\[
\sigma(u)=\sigma(v)=-\sigma(t_i).
\]

Now define
\[
q:=s_{i+1}+u=s_{i-1}+v=t_i+r_i.
\]

I will show that \(\sigma(q)\neq 0\), and that \(q\) is not one of the lonely points \(p_1,\dots,p_n\). This will contradict extremality.

### Which representations can \(q\) have?

Suppose
\[
q=s_j+t
\qquad\text{for some }t\in T.
\]
Since \(t_i\) is a vertex of \(Q\), and the two edges of \(Q\) through \(t_i\) have directions \(a_i\) and \(-a_{i-1}\), convexity gives
\[
Q\subset t_i+\operatorname{cone}(a_i,-a_{i-1}).
\]
Therefore
\[
t-t_i=r_i-s_j\in \operatorname{cone}(a_i,-a_{i-1}).
\]

On the other hand, because \(P\subset s_i+\operatorname{cone}(a_i,-a_{i-1})\), we may write uniquely
\[
s_j=s_i+\alpha a_i-\beta a_{i-1}
\qquad (\alpha,\beta\ge 0).
\]
Then
\[
r_i-s_j=(1-\alpha)a_i-(1-\beta)a_{i-1}.
\]
This belongs to \(\operatorname{cone}(a_i,-a_{i-1})\) iff
\[
0\le \alpha\le 1,\qquad 0\le \beta\le 1,
\]
i.e. iff \(s_j\in K_i\).

By Step \(1\),
\[
K_i\cap S=\{s_{i-1},s_i,s_{i+1}\}.
\]
So any representation of \(q\) must use one of the three vertices
\[
s_{i-1},\ s_i,\ s_{i+1}.
\]

Thus the only possible contributions to \(\sigma(q)\) are:

- from
  \[
  q=s_{i-1}+v,
  \]
  contributing \(\sigma(v)=-\sigma(t_i)\);

- from
  \[
  q=s_{i+1}+u,
  \]
  contributing \(\sigma(u)=-\sigma(t_i)\);

- possibly from
  \[
  q=s_i+w,
  \qquad
  w:=t_i+a_i-a_{i-1},
  \]
  if \(w\in T\), contributing \(\sigma(w)\in\{+1,-1\}\).

Therefore
\[
\sigma(q)= -2\sigma(t_i)
\quad\text{or}\quad
\sigma(q)= -2\sigma(t_i)+\sigma(w).
\]
Either way,
\[
\sigma(q)\neq 0.
\]

---

## Step \(3\): \(q\) is not one of the \(n\) lonely points

Assume
\[
q=p_k=s_k+t_k
\]
for some \(k\). Then \(q=s_k+t_k\) is a representation of the form treated above, so \(s_k\in K_i\). Hence
\[
k\in\{i-1,i,i+1\}.
\]

But each possibility is impossible:

- If \(k=i\), then
  \[
  t_i=t_i+a_i-a_{i-1},
  \]
  absurd.

- If \(k=i-1\), then
  \[
  t_{i-1}=t_i+a_i.
  \]
  Yet also
  \[
  t_i=t_{i-1}+m_{i-1}a_{i-1},
  \]
  so
  \[
  t_{i-1}=t_i-m_{i-1}a_{i-1}.
  \]
  Since \(a_i\) and \(a_{i-1}\) are nonparallel adjacent edge vectors, this is impossible.

- If \(k=i+1\), similarly
  \[
  t_{i+1}=t_i-a_{i-1}
  \]
  contradicts
  \[
  t_{i+1}=t_i+m_i a_i.
  \]

So
\[
q\notin\{p_1,\dots,p_n\}.
\]

But we have proved
\[
\sigma(q)\neq 0.
\]
Thus there is a nonzero point outside the \(n\) lonely points, contradicting
\[
\bigl|\{p:\sigma(p)\neq 0\}\bigr|=n.
\]

This contradiction proves the statement.

## Conclusion

If \(P=\operatorname{conv}(S)\) has at least three edge-direction classes, then the extremal equality
\[
\bigl|\{p:\sigma(p)\neq 0\}\bigr|=n
\]
is impossible.

Hence:

- every triangle is impossible;
- every convex \(n\)-gon with \(n\ge 5\) is impossible;
- every non-parallelogram convex quadrilateral is impossible.

Therefore the only surviving case is
\[
n=4
\quad\text{and}\quad
P\text{ is a parallelogram.}
\]

If you want, I can also write this up in a shorter “paper style” proof with the two claims compressed into a single argument.